% Source: https://artofproblemsolving.com/community/c3h2501938_two_magical_problems
% Stable identity: aops-topic-2501938-q01
We use induction on $n$.

\emph{Base case $n=1$.} A $2\times2$ board with one square removed is itself an L-tromino.

\emph{Inductive step.} Suppose every $2^n\times2^n$ board with one square removed can be tiled, and take a $2^{n+1}\times2^{n+1}$ board with one square removed. Cut it into four $2^n\times2^n$ quadrants. The removed square lies in one of them. Place one L-tromino on the three centre squares that belong to the other three quadrants (one square in each), as in the sketch for $n=2$ (the removed square is black and the central tromino is shaded).

\begin{asy}
usepackage("amsmath");
size(4.5cm);
defaultpen(fontsize(10pt));
for(int i=0;i<=8;++i){draw((i,0)--(i,8),gray(0.6));draw((0,i)--(8,i),gray(0.6));}
fill((1,6)--(2,6)--(2,7)--(1,7)--cycle,black);
fill((4,4)--(5,4)--(5,3)--(3,3)--(3,4)--(4,4)--cycle,mediumgray);
fill((4,4)--(5,4)--(5,5)--(4,5)--cycle,mediumgray);
fill((3,3)--(5,3)--(5,4)--(3,4)--cycle,mediumgray);
draw((4,3)--(5,3)--(5,5)--(4,5)--(4,4)--(3,4)--(3,3)--(4,3),linewidth(1));
draw((0,0)--(8,0)--(8,8)--(0,8)--cycle,linewidth(1.2));
draw((4,0)--(4,8),linewidth(1.2));draw((0,4)--(8,4),linewidth(1.2));
\end{asy}

Now each of the four quadrants is a $2^n\times2^n$ board with exactly one square removed or already covered. By the induction hypothesis each can be tiled, and together with the central tromino this tiles the whole board.

By induction the statement holds for all $n\ge1$. $\blacksquare$ (A count confirms that this is possible: $4^n-1$ is divisible by $3$.)
